% !PRINTED 603
\weberpage{603}
From this we obtain a basis of $\mathfrak{o}$:

a) \quad $1, \sqrt{d},$ \hfill $d\equiv 2$ or $\equiv 3\ (\mathrm{mod}\ 4),$

b) \quad $1, \dfrac{1+\sqrt{d}}{2},$ \hfill $d\equiv 1\ (\mathrm{mod}\ 4).$

If we therefore set

\weq{6}{\begin{aligned}
\text{a)}\quad \Theta &= \sqrt{d}, & d &\equiv 2, 3\ (\mathrm{mod}\ 4)\\
\text{b)}\quad \Theta &= \frac{1+\sqrt{d}}{2}, & d &\equiv 1\ (\mathrm{mod}\ 4)
\end{aligned}}
then $1, \Theta$ is a basis of $\mathfrak{o}$.

For the fundamental number we accordingly obtain, if $\Theta'$ is conjugate to $\Theta$,
\[ \Delta = \left|\begin{array}{cc} 1, & \Theta \\ 1, & \Theta' \end{array}\right|^2 = (\Theta'-\Theta)^2, \]
hence

\weq{7}{\begin{aligned}
\text{a)}\quad \Delta &= 4d, & d &\equiv 2, 3\ (\mathrm{mod}\ 4)\\
\text{b)}\quad \Delta &= d, & d &\equiv 1\ (\mathrm{mod}\ 4),
\end{aligned}}
\weq{8}{\Theta - \Theta' = \sqrt{\Delta}.}

The fundamental number is therefore $\equiv 0$ or $\equiv 1\ (\mathrm{mod}\ 4)$, and, if it is divisible by 4, must be $\equiv 8, 12\ (\mathrm{mod}\ 16)$\footnote{An integral rational number $\Delta$ that is congruent to 0 or to 1 with respect to the modulus 4 has been called by \gesperrt{Kronecker} a number of discriminant form. If $\Delta$ has the property that no square factor can be separated off from it in such a way that a number of discriminant form remains, then $\Delta$ is called a \gesperrt{stem discriminant}. The fundamental number of a quadratic field is therefore always a \gesperrt{stem discriminant}, and conversely every stem discriminant is also the fundamental number of a quadratic field. Cf. \gesperrt{H. Weber}, Zahlentheoretische Untersuchungen aus dem Gebiete der elliptischen Functionen. I, II, III. Nachrichten der Gesellschaft der Wissenschaften zu G\"ottingen. 1893.}. It can therefore (in agreement with the theorem of \gesperrt{Minkowski}) be neither $=\pm 1$ nor $=\pm 2$. The fundamental number with the smallest absolute value is $\Delta = -3.$

\begin{center}
\S.\,165.\\
\gesperrt{The Prime Ideals in $\Omega$.}
\end{center}

Now let $\varphi$ denote an arbitrary integral functional in $\Omega$. Following \S.\,146, we wish to find a basis $\alpha_1, \alpha_2$ for $\varphi$, and with it the
%CONT>>
